\medskip

\begin{enumerate}
\item 
	\begin{enumerate}
		\item L'équation paramétrique de $d_2$ montre qu'elle contient le point de coordonnées $(-3~;~0~;~5)$ et de vecteur directeur dont les composantes sont les coefficients de $k$ donc $\vect{v}\begin{pmatrix}2\\1\\0\end{pmatrix}$.
		\item Les vecteurs directeurs $\vect{u}$ et $\vect{v}$ ne sont manifestement pas colinéaires, donc les droites $d_1$ et $d_2$ ne sont pas parallèles.
		\item Une représentation paramétrique de la droite $d_1$ est $\left\{\begin{array}{l c l}
x&=2 + t\\
y&=&3 - t\\
z&=&t
\end{array}\right. \quad , t \in \R$.

S'il existe un point commun aux deux droites, il doit donc exister des réels $k$ et $t$ tels que :

$\left\{\begin{array}{l c l c l}
x &=& 2 + t&=&2k - 3\\
y&=&3 - t&=&k\\
z&=&t&=&5
\end{array}\right.$

La dernière équation donne $t = 5$, puis la deuxième $k = 3 - 5 = - 2$ et en remplaçant dans la première $2 + 5 = - 4 - 3$ : cette égalité est fausse, donc $d_1$ et $d_2$  ne sont pas sécantes.
		\item Les droites n'étant ni sécantes ni parallèles elles ne sont pas coplanaires.
	\end{enumerate}
\item 
	\begin{enumerate}
		\item On a $\vect{w}\ \cdot \vect{u} = - 1 - 2 + 3 = 0$ ;
		
		de même $\vect{w}\ \cdot \vect{v} = - 2 + 2 + 0 = 0$, donc le vecteur $\vect{w}$ est orthogonal aux vecteurs $\vect{u}$ et $\vect{v}$.
		\item Un point commun au plan $P$ et à la droite $d_2$ a ses coordonnées qui vérifient le système :
		
$\left\{\begin{array}{l c l}
x &=&2k - 3\\
y &=&k\\
z &=&5\\
5x + 4y - z - 22 &=& 0
\end{array}\right. \quad , t \in \R$, d'où en remplaçant $x,\:y,\:z$ dans la dernière équation :

$5(2k - 3) + 4(k) - 5 - 22 = 0\iff 10k - 15 + 4k - 27 = 0\iff 14k = 42 \iff 7k = 21 \iff k = 3$.

Le point commun au plan $P$ et à la droite $d_2$ a pour coordonnées (6 - 3~;~3~;~5) soit M(3~;~3~;~5).
	\end{enumerate}
\item 	
	\begin{enumerate}
		\item $\Delta$ et $d_1$ ont leurs vecteurs directeurs $\vect{w}$ et $\vect{u}$  orthogonaux donc ces droites sont orthogonales.
		
Ces deux droites sont sécantes s'il existent des réels $t$ et $r$ tels que 

$\left\{\begin{array}{l c l c l}
x&=&2 + t&=&- r + 3\\
y&=&3 - t&=&2r + 3\\
z&=&t   &=&3r + 5
\end{array}\right. \quad , t, r  \in \R$.

La dernière équation $t = 3r + 5$ donne en remplaçant dans la deuxième : 

$3 - 3r - 5 = 2r + 3 \iff - 5 = 5r \iff r = - 1$, d'où $t = 5 - 3 = 2$ et en remplaçant dans la première équation on obtient $2 + 2 = 1 + 3$, égalité vraie.

$\Delta$ et $d_1$ sont sécantes au point L(4~;~1~;~2)
		\item  Conclusion : on a trouvé une droite $\Delta$ perpendiculaire commune aux deux droites $d_1$ et $d_2$.
	\end{enumerate}
\end{enumerate}

\bigskip
